
\subsection{Proof of Theorem~\ref{thm:dual-iqc}}\label{proof:dual-iqc}
$1)\implies2)$ Assume there exists a $\rho<1$ such that
$\|p_T G_0\tilde w\|_{L_{[a,b]}}\leq\rho\|p_T \tilde w\|_{L_{[a,b]}}$
for all $\tilde{w}\in L_{[a,b]}$ and $T\geq0$. We need to show that
$\mcl I\leq-\epsilon\|p_Tw\|_{L_{[a,b]}}^2$ for some $\epsilon>0$. First since $\Gamma:=\Theta_{21}P+\Theta_{22}$ is causally boundedly invertible,
$\Theta\bmat{P\\I}w=\bmat{G_0\\I}\Gamma w$, and hence
$\mcl I=\|p_TG_0\Gamma  w\|_{L_{[a,b]}}^2-\|p_T\Gamma  w\|_{L_{[a,b]}}^2$. By statement 1 and causality of $G_0$,
\[
    \mcl I\leq-(1-\rho^2)\|p_T\Gamma w\|_{L_{[a,b]}}^2
    \leq-\frac{1-\rho^2}{\|\Gamma^{-1}\|^2}\|p_Tw\|_{L_{[a,b]}}^2,
\]
where the second inequality follows from
$\|p_Tw\|_{L_{[a,b]}}=\|p_T\Gamma^{-1}p_T\Gamma w\|_{L_{[a,b]}}
\leq\|\Gamma^{-1}\|\|p_T\Gamma w\|_{L_{[a,b]}}$. Thus, $\epsilon=\frac{(1-\rho^2)}{\|\Gamma^{-1}\|^2}>0$.

$2)\implies1)$ Assume that there exists a $\epsilon > 0$ such that $\mcl I\leq-\epsilon\|p_Tw\|_{L_{[a,b]}}^2$ for every $w$, and every $T\geq0$. 
We need to establish there exists a $\rho<1$ such that
$\|p_T G_0\tilde w\|_{L_{[a,b]}}\leq\rho\|p_T \tilde w\|_{L_{[a,b]}}$
for all $\tilde{w}\in L_{[a,b]}$ and $T\geq0$. For arbitrary $\tilde w$, 
set $w=\Gamma^{-1}\tilde w$. Then
\[
    \|p_TG_0\tilde w\|_{L_{[a,b]}}^2-\|p_T\tilde w\|_{L_{[a,b]}}^2
    \leq-\epsilon\|p_T\Gamma^{-1}\tilde w\|_{L_{[a,b]}}^2.
\]
Since $\Gamma$ is bounded and causal,
$\|p_T\tilde w\|_{L_{[a,b]}}\leq
\|\Gamma\|\|p_T\Gamma^{-1}\tilde w\|_{L_{[a,b]}}$. Therefore,
\[
    \|p_TG_0\tilde w\|_{L_{[a,b]}}^2
    \leq\left(1-\frac{\epsilon}{\|\Gamma\|^2}\right)
    \|p_T\tilde w\|_{L_{[a,b]}}^2.
\]
Choose $\rho<1$ such that $1-\epsilon/\|\Gamma\|^2\leq\rho^2<1$, 
then $\|p_TG_0\tilde w\|_{L_{[a,b]}}\leq\rho\|p_T\tilde w\|_{L_{[a,b]}}$
for every $\tilde w$ and $T\geq0$.

$1)\iff3)$ We show that there exists a causally boundedly invertible
$\underline\Gamma$ such that
\[
    \mbf D(\Theta)\bmat{P^\top\\I}\underline w
    =\bmat{G_0^\top\\I}\underline\Gamma\underline w.
\]
then we can apply the previous arguments to show $1)\iff 3)$. Define
\[
    S:=\Theta_{11}-\Theta_{12}\Theta_{22}^{-1}\Theta_{21}.
\]
Since $\Theta$ and $\Theta_{22}$ are boundedly invertible, so is $S$ by Schur complement.
Block inversion of $\mbf D(\Theta)$ gives
\[
    \mbf D(\Theta)_{21}
    =(S^\top)^{-1}\Theta_{21}^\top(\Theta_{22}^\top)^{-1},
    \qquad
    \mbf D(\Theta)_{22}=(S^\top)^{-1}.
\]
Therefore, define
\[
    \underline\Gamma
    :=\mbf D(\Theta)_{21}P^\top+\mbf D(\Theta)_{22}
    =(S^\top)^{-1}
      \bigl(I+P\Theta_{22}^{-1}\Theta_{21}\bigr)^\top.
\]
Now $\Gamma=\Theta_{22}(I+\Theta_{22}^{-1}\Theta_{21}P)$.
Thus, invertibility of $\Gamma$ and the push-through rule imply that
$I+P\Theta_{22}^{-1}\Theta_{21}$, and hence $\underline\Gamma$, is
causally boundedly invertible. Normalizing the upper block gives
\[
    \mbf D(\Theta)\bmat{P^\top\\I}\underline w
    =\bmat{\underline G_0\\I}\underline\Gamma\underline w,
    \qquad
    \underline G_0:=(\Theta_{22}^\top)^{-1}
    (P^\top\underline\Gamma^{-1}+\Theta_{12}^\top).
\]
Using
\[
    \underline\Gamma^{-\top}
    =S(I+P\Theta_{22}^{-1}\Theta_{21})^{-1},
    \qquad
    (I+P\Theta_{22}^{-1}\Theta_{21})^{-1}P\Theta_{22}^{-1}
    =P\Gamma^{-1},
\]
where the second identity is the push-through rule, yields
\begin{align*}
    \underline G_0
    &=(\Theta_{22}^\top)^{-1}
      (P^\top\underline\Gamma^{-1}+\Theta_{12}^\top) \\
    &=\left(
        \underline\Gamma^{-\top}P\Theta_{22}^{-1}
        +\Theta_{12}\Theta_{22}^{-1}
      \right)^\top \\
    &=\left[
        S(I+P\Theta_{22}^{-1}\Theta_{21})^{-1}
        P\Theta_{22}^{-1}
        +\Theta_{12}\Theta_{22}^{-1}
      \right]^\top \\
    &=\left[
        SP\Gamma^{-1}+\Theta_{12}\Theta_{22}^{-1}
      \right]^\top \\
    &=\left[
        \left(SP+\Theta_{12}\Theta_{22}^{-1}\Gamma\right)
        \Gamma^{-1}
      \right]^\top.
\end{align*}
We now show that the first part reduces to
\begin{align*}
    SP+\Theta_{12}\Theta_{22}^{-1}\Gamma
    &=\Theta_{11}P
      -\Theta_{12}\Theta_{22}^{-1}\Theta_{21}P \\
    &\quad+\Theta_{12}\Theta_{22}^{-1}\Theta_{21}P
      +\Theta_{12} \\
    &=\Theta_{11}P+\Theta_{12}.
\end{align*}
Consequently,
\[
    \underline G_0
    =\left[(\Theta_{11}P+\Theta_{12})\Gamma^{-1}\right]^\top
    =G_0^\top,
\]
and hence
\[
    \mbf D(\Theta)\bmat{P^\top\\I}\underline w
    =\bmat{G_0^\top\\I}\underline\Gamma\underline w.
\]
By Theorem~\ref{thm:dual-l2-gain}, $G_0$ and $G_0^\top$ have the
same induced $L_2$-gain. Hence, for any $\rho$,
\[
\begin{aligned}
&\|p_TG_0\tilde w\|_{L_{[a,b]}}
 \leq\rho\|p_T\tilde w\|_{L_{[a,b]}}
 \quad\text{for all $\tilde w$ and $T\geq0$}\\
&\qquad\Longleftrightarrow\qquad
\|p_TG_0^\top\underline{\tilde w}\|_{L_{[a,b]}}
 \leq\rho\|p_T\underline{\tilde w}\|_{L_{[a,b]}}
 \quad\text{for all $\underline{\tilde w}$ and $T\geq0$}.
\end{aligned}
\] Since
$\underline\Gamma$ is causally boundedly invertible, the same arguments used
for $1)\implies2)$ and $2)\implies1)$, with
$G_0$, $\Gamma$, $w$, and $\mcl I$ replaced by
$G_0^\top$, $\underline\Gamma$, $\underline w$, and
$\underline{\mcl I}$, respectively, establish $1)\iff3)$.

\paragraph{Realization of $G_0$.}
Let $\mbf x$ collect the plant, filter, and any controller states, so that the
controller is already included in the primal operators. Partition the filtered realization as
\[
\bmat{\mcl T\dot{\mbf x}\\\tilde z\\\tilde w}
=
\bmat{\mcl A&\mcl B\\\mcl C_{\mrm{z}}&\mcl D_{\mrm{z}}\\
\mcl C_{\mrm{w}}&\mcl D_{\mrm{w}}}
\bmat{\mbf x\\w}.
\]
If $\mcl D_{\mrm{w}}$ is invertible, eliminating $w$ gives the realization
\begin{equation}\label{eq:G0-primal-realization}
G_0\sim
\left\{
\mcl T,
\mcl A-\mcl B\mcl D_{\mrm{w}}^{-1}\mcl C_{\mrm{w}},
\mcl B\mcl D_{\mrm{w}}^{-1},
\mcl C_{\mrm{z}}-\mcl D_{\mrm{z}}\mcl D_{\mrm{w}}^{-1}\mcl C_{\mrm{w}},
\mcl D_{\mrm{z}}\mcl D_{\mrm{w}}^{-1}
\right\},
\end{equation}
with input $\tilde w$ and output $\tilde z$.
